\section{Summary and Outlook}
\label{sec:summary}

In this work, we have calculated the CS kernel on a MILC lattice configuration in the large momentum effective theory framework. Comparing with our previous studied~\cite{LatticeParton:2020uhz}, the one-loop matching kernel has been adopted in this study, and several hadron momenta were used to extract the CS kernel. We found that in the small $b_{\perp}$ region, our results are consistent with perturbative QCD. 
In large $b_{\perp}$ region, our results seem consistent with other lattice calculations in the literature within uncertainties. 

For our future studies, we need to
use lattice configurations with multiple lattice spacing to understand the finite
lattice spacing effects. We would
use a valence quark mass consistent with the sea quark one to reduce the non-unitarity effects. One such effect 
might be the imaginary part of the meson
wave function which seems inconsistent
with perturbative calculation at present time. Clearly, all these explorations will
take more computational resources. 
